\begingroup
\section{Node QB: quadratic-bubble reduction}
\label{module:QB}
\noindent\textit{Audited source: \nolinkurl{QB_bubble_reduction.tex}.}
\subsection{Quadratic polar spline}

Let
\begin{equation}\label{QB:eq:fan}
 \alpha_i>0,\qquad \sum_i\alpha_i=2\pi,
 \qquad a=\frac{2\pi}{N},\qquad S=\max_i\alpha_i+a.
\end{equation}
Choose normal angles $\theta_i$ with
$\theta_{i+1}-\theta_i=\alpha_i$ and put
\begin{equation}\label{QB:eq:vertex}
 \varphi_i=\theta_i+\frac{\alpha_i}{2}.
\end{equation}
The polar cell of side $i$ is
\begin{equation}\label{QB:eq:cell}
 K_i=[\varphi_{i-1},\varphi_i]
 =[\theta_i-\alpha_{i-1}/2,\theta_i+\alpha_i/2].
\end{equation}
Define the continuous periodic quadratic spline
\begin{equation}\label{QB:eq:q}
 q_\alpha(\theta)=\frac12(\theta-\theta_i)^2,
 \qquad \theta\in K_i.
\end{equation}
Continuity follows because the two values at $\varphi_i$ are both
$\alpha_i^2/8$.

The exact radial fan graph is
\begin{equation}\label{QB:eq:f}
 f_\alpha(\theta)=h_\alpha\sec(\theta-\theta_i)-1,
 \qquad \theta\in K_i,
\end{equation}
with
$h_\alpha=(\pi/\sum_j\tan(\alpha_j/2))^{1/2}$.

\begin{lemma}[Secant/quadratic profile split]
\label{QB:lem:split}
There is a constant $c_\alpha$ and a continuous remainder $r_\alpha$ such
that
\begin{equation}\label{QB:eq:split}
 f_\alpha=c_\alpha+q_\alpha+r_\alpha,
\end{equation}
and, without a minimum-angle assumption,
\begin{equation}\label{QB:eq:r-size}
 \|r_\alpha\|_\infty\le CS^4,
 \qquad \|r_\alpha\|_{W^{1,\infty}}\le CS^3,
 \qquad \|r_\alpha\|_{H^{1/2}}^2\le CS^7.
\end{equation}
Also
\begin{equation}\label{QB:eq:q-size}
 \|q_\alpha\|_\infty\le CS^2,
 \qquad \|q_\alpha\|_{W^{1,\infty}}\le CS,
 \qquad \|q_\alpha\|_{H^{1/2}}^2\le CS^3.
\end{equation}
\end{lemma}

\begin{proof}
Take $c_\alpha=h_\alpha-1$ and write, on a cell,
\begin{equation*}
 r_\alpha(x)
 =(h_\alpha-1)\frac{x^2}{2}
 +h_\alpha\left(\sec x-1-\frac{x^2}{2}\right).
\end{equation*}
Since $h_\alpha-1=O(\sum_i\alpha_i^3)=O(S^2)$ and
$\sec x-1-x^2/2=O(x^4)$ through one normalized derivative, the first two
estimates in \eqref{QB:eq:r-size} follow.  Same and adjacent half-cells in the
Slobodeckij seminorm cost $C\alpha_i^8$.  On separated half-cells subtract
both means and use the first-separating dyadic annulus.  Thus
\begin{equation*}
 \|r_\alpha\|_{H^{1/2}}^2
 \le C\sum_i\alpha_i^8
 \le CS^4\sum_i\alpha_i^4\le CS^7.
\end{equation*}
The proof of \eqref{QB:eq:q-size} is identical with local amplitude
$\alpha_i^2$ and derivative amplitude $\alpha_i$.
\end{proof}

\subsection{The higher secant profile is harmless}

Recall the symmetric null form of stage 138,
\begin{multline}\label{QB:eq:C}
 \Cnull(g_1,g_2,g_3)\\
 =\frac13\sum_{\rm cyc}\int_\T g_1
 \{(\Dop g_2)(\Dop g_3)-g_2'g_3'\}\,d\theta.
\end{multline}
It has total order two and the tame estimate
\begin{equation}\label{QB:eq:C-tame}
 |\Cnull(g_1,g_2,g_3)|
 \le C\sum_{p<q}\|g_p\|_{H^{1/2}}\|g_q\|_{H^{1/2}}
 \|g_\ell\|_{W^{1,\infty}}.
\end{equation}

\begin{lemma}[Constants disappear]
\label{QB:lem:constant}
For every $g$ and every constant $c$,
\begin{equation}\label{QB:eq:constant}
 \Cnull(g+c,g+c,g+c)=\Cnull(g,g,g).
\end{equation}
The same is true after polarization whenever a constant occupies any one
slot.
\end{lemma}

\begin{proof}
Terms in which a differentiated slot is constant vanish.  If the
undifferentiated slot is constant, Parseval gives
$\|\Dop g\|_2^2=\|g'\|_2^2$.  Polarization proves the last assertion.
\end{proof}

\begin{theorem}[Secant-to-quadratic transfer]
\label{QB:thm:transfer}
There are moduli $\varepsilon_q(S),\varepsilon_{q,s}(S)\to0$ such that
\begin{align}
 &|\{\Cnull(f_\alpha^3)-\Cnull(q_\alpha^3)\}
 -\{\Cnull(f_{a\mathbf1}^3)-\Cnull(q_{a\mathbf1}^3)\}|
 \le\varepsilon_q(S)\Jfour(\alpha),
 \label{QB:eq:transfer-fan}\\
 &|\Cnull(f_\alpha,f_\alpha,v_\alpha(z))
 -\Cnull(q_\alpha,q_\alpha,v_\alpha(z))|
 \le\varepsilon_{q,s}(S)\Jfour(\alpha)^{1/2}
 \widetilde{\mathfrak G}_\alpha(z)^{1/2}.
 \label{QB:eq:transfer-mixed}
\end{align}
Here $v_\alpha(z)$ is the normal component of the physical support trace.
\end{theorem}

\begin{proof}
By Lemma~\ref{QB:lem:constant}, only $r_\alpha$ occurs in the difference.
Equations \eqref{QB:eq:C-tame}--\eqref{QB:eq:q-size} give the absolute estimates
\begin{equation}\label{QB:eq:transfer-absolute}
 |\Cnull(f_\alpha^3)-\Cnull(q_\alpha^3)|\le CS^6,
\end{equation}
and
\begin{equation}\label{QB:eq:transfer-mixed-absolute}
 |\Cnull(f_\alpha,f_\alpha,v)-\Cnull(q_\alpha,q_\alpha,v)|
 \le CS^{9/2}\|v\|_{H^{1/2}}.
\end{equation}
For example, an energetic $(q,r)$ pair and a coefficient $q$ cost
$S^{3/2}S^{7/2}S=S^6$; the other placements are no larger.

Use the far threshold $\Jfour\ge S^{6-\beta}$, $0<\beta<1$.
Equations \eqref{QB:eq:transfer-absolute}--
\eqref{QB:eq:transfer-mixed-absolute} give an $o(1)$ multiple of the required
right sides.  In the complementary regime,
\begin{equation*}
 \|\alpha-a\mathbf1\|_\infty/S
 \le CS^{(2-\beta)/4},
\end{equation*}
so the whole segment is in a fixed complex material polydisc.  Pull every
half-cell to $[0,1]$ before differentiating.  Cauchy's estimate applied to
the two absolute majorants gives
\begin{equation*}
 |F_q''(t)|\le CS^2E(t),
 \qquad \|\Psi_q'(t)\|\le CS^{5/2}E(t)^{1/2}.
\end{equation*}
Cyclic symmetry kills the regular first fan derivative and the regular
support functional on the exact quotient.  Integrate with the weighted and
unweighted quartic actions.  Physical trace equivalence replaces
$\|v_\alpha(z)\|_{H^{1/2}}$ by
$C\widetilde{\mathfrak G}_\alpha(z)^{1/2}$.
\end{proof}

\subsection{Exact midpoint-atom formula}

Let
\begin{equation}\label{QB:eq:mu}
 \mu_\alpha=\sum_i\alpha_i\delta_{\varphi_i}.
\end{equation}

\begin{lemma}[Distributional spline identity]
\label{QB:lem:distribution}
One has
\begin{equation}\label{QB:eq:distribution}
 q_\alpha''=1-\mu_\alpha
\end{equation}
on $\T$.  With Fourier convention
$\widehat g(n)=(2\pi)^{-1}\int_\T g(\theta)e^{-in\theta}\,d\theta$, put
\begin{equation}\label{QB:eq:M}
 M_\alpha(n)=\sum_i\alpha_i e^{-in\varphi_i}.
\end{equation}
Then, for $n\ne0$,
\begin{equation}\label{QB:eq:qhat}
 \widehat q_\alpha(n)=\frac{M_\alpha(n)}{2\pi n^2}.
\end{equation}
Moreover the midpoint cancellation is the exact identity
\begin{equation}\label{QB:eq:midpoint-cancellation}
 M_\alpha(n)=\sum_i\alpha_i e^{-in\varphi_i}
 \left\{1-\operatorname{sinc}\left(\frac{n\alpha_i}{2}\right)\right\},
 \qquad n\ne0.
\end{equation}
\end{lemma}

\begin{proof}
Inside every open cell, $q_\alpha''=1$.  At the vertex $\varphi_i$ the
one-sided derivatives are $\alpha_i/2$ and $-\alpha_i/2$, so the derivative
jump is $-\alpha_i$.  This proves \eqref{QB:eq:distribution}; Fourier
transformation gives \eqref{QB:eq:qhat}.

The normal intervals
$I_i=[\theta_i,\theta_{i+1}]$ partition the circle and have midpoint
$\varphi_i$.  Hence, for $n\ne0$,
\begin{equation*}
 0=\int_\T e^{-in\theta}\,d\theta
 =\sum_i\alpha_i e^{-in\varphi_i}
 \operatorname{sinc}\left(\frac{n\alpha_i}{2}\right).
\end{equation*}
Subtract this equality from \eqref{QB:eq:M}.
\end{proof}

\begin{theorem}[Exact discrete cubic scalar]
\label{QB:thm:discrete}
The pure quadratic-bubble null form is
\begin{equation}\label{QB:eq:discrete}
 \boxed{
 \Cnull(q_\alpha,q_\alpha,q_\alpha)
 =\frac1{\pi^2}\operatorname{Re}
 \sum_{p,q\ge1}
 \frac{M_\alpha(p)M_\alpha(q)
 \overline{M_\alpha(p+q)}}{pq(p+q)^2}.}
\end{equation}
The series is interpreted by symmetric Abel limits and is absolutely
convergent after the midpoint cancellation \eqref{QB:eq:midpoint-cancellation}
is inserted sourcewise.
\end{theorem}

\begin{proof}
For a real $g$, Fourier expansion of \eqref{QB:eq:C} shows that only triples
with two frequencies of the same sign survive.  If $p,q>0$ and the third
frequency is $-(p+q)$, the unsymmetrized coefficient is $2pq$.
The positive and negative triples are conjugate, so
\begin{equation*}
 \Cnull(g,g,g)=8\pi\operatorname{Re}
 \sum_{p,q\ge1}pq\,\widehat g(p)\widehat g(q)
 \overline{\widehat g(p+q)}.
\end{equation*}
Insert \eqref{QB:eq:qhat}; since
$8\pi/(2\pi)^3=1/\pi^2$, this gives \eqref{QB:eq:discrete}.
Abel regularization $r^{p+q}$ justifies rearrangement.  Formula
\eqref{QB:eq:midpoint-cancellation} gives two powers of $n\alpha_i$ at low
frequency and the trivial bound at high frequency; dyadic source
allocation then makes the Abel limit independent of the order of
summation.
\end{proof}

\subsection{Final discrete obligations}

\begin{obligation}[DF: midpoint cubic Bregman estimate]
\label{QB:obl:DF}
For the explicit scalar on the right of \eqref{QB:eq:discrete}, prove
\begin{equation}\label{QB:eq:DF}
 |\Cnull(q_\alpha^3)-\Cnull(q_{a\mathbf1}^3)|
 \le CS\Jfour(\alpha).
\end{equation}
Every cumulative phase produced by differentiating $\varphi_i$ must first
be summed cyclically; the midpoint factor in
\eqref{QB:eq:midpoint-cancellation} is not to be differentiated row by row.
\end{obligation}

\begin{obligation}[DM: polarized midpoint/support estimate]
\label{QB:obl:DM}
For the polarization of \eqref{QB:eq:discrete} with one slot replaced by the
physical normal trace $v_\alpha(z)$, prove
\begin{equation}\label{QB:eq:DM}
 |\Cnull(q_\alpha,q_\alpha,v_\alpha(z))|
 \le\varepsilon_{\rm DM}(S)\Jfour(\alpha)^{1/2}
 \widetilde{\mathfrak G}_\alpha(z)^{1/2},
 \qquad \varepsilon_{\rm DM}(S)\to0.
\end{equation}
\end{obligation}

\begin{proposition}[Equivalence]
\label{QB:prop:equivalence}
DF and the secant-to-quadratic transfer imply the pure principal row NF.
DF, DM, and the same transfer imply the mixed principal row NM.  These are
exactly the two terminal principal inputs exported by QB; the conversion
to CF, CM, NR, and J0 is performed only in the current
\texttt{NR\_assembly.tex} node.
\end{proposition}

\begin{proof}
Apply Theorem~\ref{QB:thm:transfer} to the pure and mixed null forms and absorb
its moduli into the DF and DM moduli.  No downstream tensor or value
statement is imported here.
\end{proof}
\endgroup
